\chapter{Généralités sur les fonctions}
\thispagestyle{fancy}


\section{Définitions et notations}

\subsection{Définition}

\begin{bexemple}
	On considère la fonction \(f\) qui exprime l'aire d'un rectangle de dimensions 3 et \(x\).\\
	Une expression littérale de \(f\) est donc : \(f(x)=3 x\).
\end{bexemple}
\begin{bdefi}
	Une fonction \(f\) associe à tout nombre réel \(x\) un unique nombre réel, noté \(f(x)\).\\
	On note également : \(x \mapsto f(x)\) ou \(y=f(x)\).
\end{bdefi}


\subsection{Image et antécédent}

\begin{bexemple}
	Dire que : \(f(2)=5\) signifie que : 
	
	\includegraphics[scale=0.25]{Cours_img01}
	
	On dit que :
	\begin{itemEx}
		\item l'image de 2 par la fonction \(f\) est 5.
		\item un antécédent de 5 par \(f\) est 2.
	\end{itemEx}
\end{bexemple}

\begin{brmqs}
	\begin{itemRem}
		\item Un nombre possède une unique image.
		\item Cependant, un nombre peut posséder plusieurs antécédents.
	\end{itemRem}
\end{brmqs}

\begin{bmethodet}{Déterminer l'image d'une fonction par calcul}
	\grasInMet{Enoncé:}
	
	Soit la fonction \(g\) définie par \(g(x)=x^{2}-2\).\\
	Calculer l'image de 6 par la fonction \(g\).
	
	\grasInMet{Réponse:}
	$$g(x)=x^{2}-2 \iff g(6)=6^{2}-2 \iff g(6)=36-2 \iff \fbox{g(6)=34}$$
	L'image de 6 par la fonction \(g\) est 34.
\end{bmethodet}


\begin{bmethodet}{Déterminer un antécédent par calcul}
	\grasInMet{Enoncé:}
	
	Soit la fonction \(f\) définie par \(f(x)=2 x-3\).\\
	Déterminer un antécédent de $-5$ par la fonction \(f\).
	
	\grasInMet{Réponse:}
	
	On cherche un antécédent de $-5$ donc $-5$ est une image.\\
	On peut donc écrire : \(f(x)=-5\)\\
	Soit : \(2 x-3=-5\)\\
	On résout ainsi l'équation :
	$$2 x=3-5 \iff 2 x=-2 \iff \fbox{x=-1}$$
	L'antécédent de $-5$ par \(f\) est donc $-1$.
\end{bmethodet}

\section{Représentation graphique}

\begin{bmethodet}{Représenter graphiquement une fonction}
	\grasInMet{Enoncé:}
	
	Soit la fonction \(f\) définie par \(f(x)=5 x-x^{2}\).\\
	On donne un tableau de valeurs de la fonction \(f\) :
	\begin{center}
		\begin{tabular}[t]{|l|l|l|l|l|l|l|l|l|}
			\hline
			\(x\) & 1 & 1,5 & 2 & 2,5 & 3 & 3,5 & 4 & 4,5 \\
			\hline
			\(f(x)\) & 4 & 5,25 & 6 & 6,25 & 6 & 5,25 & 4 & 2,25 \\
			\hline
		\end{tabular}
	\end{center}
	Tracer, dans un repère, la courbe représentative de la fonction \(f\).
	
	\grasInMet{Réponse:}
	
	\begin{minipage}[c]{0.48\linewidth}
		On représente les données du tableau de valeurs dans un repère tel qu'on trouve en abscisse les valeurs de \(x\) et en ordonnée les valeurs de \(f(x)\) correspondantes.
		
		En reliant les points, on obtient une courbe.\\
		Tout point de la courbe possède donc des coordonnées de la forme \((x ; f(x))\).
	\end{minipage} \hfill
	\begin{minipage}[c]{0.48\linewidth}
		\includegraphics[scale=0.25]{Cours_img02}
	\end{minipage}
	
	
	
\end{bmethodet}

\section{Résolution graphique d'équations et d'inéquations}

\begin{bmethodet}{Résoudre graphiquement une équation}
	\grasInMet{Enoncé:}
	
	\begin{minipage}[c]{0.48\linewidth}
		On a représenté la courbe de la fonction \(f\) définie par \(f(x)=5 x-x^{2}\).\\
		Résoudre graphiquement l'équation \(5 x-x^{2}=4\).\\
	\end{minipage} \hfill
	\begin{minipage}[c]{0.48\linewidth}
		\includegraphics[scale=0.25]{Cours_img03}
	\end{minipage}
	
	\grasInMet{Réponse:}
	
	\begin{minipage}[c]{0.48\linewidth}
		L'équation \(5 x-x^{2}=4\) peut s'écrire \(f(x)=4\).\\
		Ce qui revient à trouver les antécédents de 4 par la fonction \(f\).\\
		On \guillemets{part} de l'ordonnée 4, on \guillemets{rejoint} la courbe et on lit les solutions sur l'axe des abscisses : \(x=1\) ou \(x=4\).
		
		On peut noter : \(S=\{1 ; 4\}\).
	\end{minipage} \hfill
	\begin{minipage}[c]{0.48\linewidth}
		\includegraphics[scale=0.25]{Cours_img04}
	\end{minipage}
	
	
\end{bmethodet}

\begin{brmqs}
	\begin{itemRem}
		\item Par lecture graphique, les solutions obtenues sont approchées.
		\item L'équation \(f(x)=7\), par exemple, ne semble pas avoir de solution car la courbe représentée ne possède pas de point d'ordonnée 7.
		\item Graphiquement, on ne peut pas être certain que les solutions qui apparaissent sont les seules. Il pourrait y en avoir d'autres au-delà des limites de la représentation graphique tracée.
	\end{itemRem}
\end{brmqs}

\begin{bmethodet}{Résoudre graphiquement une inéquation}
	\grasInMet{Enoncé:}
	
	Dans la méthode précédente, on a représenté la courbe de la fonction \(f\) définie par \(f(x)=5 x-x^{2}\).\\
	Résoudre graphiquement l'inéquation \(5 x-x^{2}>4\).
	
	\grasInMet{Réponse:}
	
	\begin{minipage}[c]{0.48\linewidth}
		L'inéquation \(5 x-x^{2}>4\) peut s'écrire \(f(x)>4\). Ce qui revient à déterminer les points de la courbe dont l'ordonnée est strictement supérieure à 4. On lit les solutions correspondantes sur l'axe des abscisses :\\
		\(x\) est strictement compris entre 1 et 4.\\
		On peut noter : \(S=] 1 ; 4[\).\\
	\end{minipage} \hfill
	\begin{minipage}[c]{0.48\linewidth}
		\includegraphics[scale=0.25]{Cours_img05}
	\end{minipage}
	
	
	
\end{bmethodet}

\section{Variations d'une fonction}

\subsection{Taux de variation}

\begin{bdefi}
	Le taux de variation de la fonction \(f\) entre \(a\) et \(b\) est le nombre: \(\dfrac{f(b)-f(a)}{b-a}\)
\end{bdefi}

\begin{bprop}
	Le taux de variation de \(f\) entre \(a\) et \(b\) est la pente de la droite passant par les points d'abscisses \(a\) et \(b\) de la courbe de \(f\).
\end{bprop}

\begin{bmethodet}{Déterminer un taux de variation d'une fonction}
	\grasInMet{Enoncé:}
	
	Soit \(f\) la fonction définie sur \(\mathbb{R}\) par : \(f(x)=2 x^{2}+1\).
	
	\begin{enumMet}
		\item Déterminer le taux de variation entre 1 et 3.
		\item Interpréter géométriquement ce taux de variation.
	\end{enumMet}
	
	\grasInMet{Réponse:}
	
	\begin{minipage}[c]{0.58\linewidth}
		\begin{enumMet}
			\item Si \(f(x)=2 x^{2}+1\), alors le taux de variation de \(f\) entre 1 et 3 est égal à :
			\[
			\frac{f(3)-f(1)}{3-1}=\frac{2 \times 3^{2}+1-\left(2 \times 1^{2}+1\right)}{2}=\frac{19-3}{2}=8
			\]
			
			\item Le taux de variation de \(f\) entre 1 et 3 est égal à 8 donc la pente de la droite passant par les points d'abscisses 1 et 3 est égale à 8.\\
		\end{enumMet}
	\end{minipage} \hfill
	\begin{minipage}[c]{0.38\linewidth}
		\includegraphics[scale=0.18]{Cours_img06}\\
	\end{minipage}
	
	
\end{bmethodet}

\subsection{Fonctions monotones}

\begin{bdefi}
	On dit qu'une fonction \(f\) est monotone sur un intervalle I, si \(f\) est :
	\begin{itemDef}
		\item soit croissante sur I,
		\item soit décroissante sur I,
		\item soit constante sur I.
	\end{itemDef}
\end{bdefi}
\newpage

\begin{bprops}
	\begin{itemProp}
		\item Si le taux de variation d'une fonction \(f\) entre deux nombres quelconques d'un intervalle I est positif, alors \(f\) est strictement croissante sur I.
		\item S'il est négatif, \(f\) est strictement décroissante sur I.
		\item S'il est nul, \(f\) est constante sur I.
	\end{itemProp}
\end{bprops}

\begin{bmethodet}{Étudier les variations d'une fonction à l'aide du taux de variation}
	\grasInMet{Enoncé:}
	
	Soit \(f\) la fonction définie sur \(\mathbb{R}\) par : \(f(x)=5 x-3\).\\
	Démontrer que \(f\) est strictement croissante sur \(\mathbb{R}\).
	
	\grasInMet{Réponse:}
	
	On considère deux nombres quelconques distincts \(a\) et \(b\).\\
	Le taux de variation de \(f\) entre \(a\) et \(b\) est égal à :
	
	\[
	\frac{f(b)-f(a)}{b-a}=\frac{5 b-3-(5 a-3)}{b-a}=\frac{5 b-5 a}{b-a}=\frac{5(b-a)}{b-a}=5
	\]
	
	Or, \(5>0\) donc \(\dfrac{f(b)-f(a)}{b-a}>0\) et donc \(f\) est strictement croissante sur \(\mathbb{R}\).	
\end{bmethodet}
